% =====================================================================================
% Figure -- the sequential workflow of the user study session.
% Belongs in \subsection{Experimental Procedure} (methodology_section.tex, line ~225),
% referenced from the sentence introducing the enumerated procedure.
%
% No extra packages required.
%
% Expected image file, relative to the main .tex:
%     figures/study-workflow.png
%
% The source render is roughly 16:9 (1568 x 900) and is mostly small text, so it needs
% the full text width -- hence figure*, not figure. If the paper is single-column, swap
% figure* for figure and the caption still applies unchanged. Export at >= 300 dpi or,
% better, as PDF/SVG from the diagram tool: the panel labels will not survive a
% low-resolution raster in print.
%
% CAPTION CAVEAT -- read before using. The diagram shows one uninterrupted run of six
% stages, but methodology_section.tex (lines 227--244) records that the session was
% physically SPLIT INTO TWO PARTS by a headset reboot between Trial 2 and Trial 3, to
% avoid the CVIP memory exhaustion crash, with participants relocating indoors and the
% markers re-scanned afterwards. Two options, pick one:
%   (a) add the break to the diagram and then uncomment the extra caption sentence below;
%   (b) leave the diagram as-is and keep the sentence, which tells the reader the figure
%       is a logical rather than a wall-clock ordering.
% Do not ship the figure with neither: as drawn, it contradicts the procedure text.
%
% A second mismatch to check: the diagram labels panel (b) "select recalled objects as
% present / absent / unsure", while the text describes a 30-item recall test with a
% 7-point confidence rating on Trial 1. Confirm the response format before submitting.
%
% Reference in text as Figure~\ref{fig:procedure}.
% =====================================================================================

\begin{figure*}[t]
  \centering
  \includegraphics[width=\textwidth]{figures/study-workflow}
  \caption{Sequential workflow of a study session. A session runs through six stages:
    informed consent, the pre-study survey capturing demographics and baseline measures,
    eye tracking calibration on the Magic Leap 2, four experimental trials, the post-study
    survey, and a closing open-ended interview. The four trials are identical in
    structure and differ only in the condition assigned by the counterbalancing matrix
    (Table~\ref{tab:counterbalance}). Each trial runs the same three steps in order: the
    participant searches the site for virtual gems while incidentally encountering recall
    objects (a); then completes the recall test, judging each object present, absent, or
    unsure (b); then completes the post-trial survey battery (c). This block repeats in
    full for all four trials before the session advances to the post-study survey.
    % Uncomment if the diagram is not redrawn to show the break -- see caveat above.
    % The figure gives the logical ordering of the session; in practice trials were split
    % across a headset reboot between Trial~2 and Trial~3, as described below.
    }
  \Description{A flowchart of a study session, read left to right. A row of six colored
    boxes across the top gives the stages in order: 1. Consent (participant onboarding),
    2. Pre-Study Survey (baseline responses), 3. Eye Calibration (Magic Leap 2 setup),
    4. Trials times four (repeated experimental sequence), 5. Post-Study Survey (overall
    reflections), and 6. Final Interview (open-ended question). Below, a larger panel
    expands each stage into a box connected by arrows in the same order: participants sign
    consent form; fill out pre-study survey; eye calibration on Magic Leap 2; then the
    trials panel; then fill out post-study survey; and finally, below that and reached by
    a downward arrow, an open-ended final interview. The trials panel is drawn as a large
    enclosing box titled Trials times four, noting that Trial 1, Trial 2, Trial 3, and
    Trial 4 follow the same sequence. Inside it, three stacked boxes are joined by
    downward arrows: a, Search and Recall, search for virtual gems and recall objects in
    the scene; b, Recall Task, select recalled objects as present, absent, or unsure; and
    c, Post-Trial Survey, fill out post-trial survey after each trial. A banner beneath
    the panel reads that the complete trial block repeats for all four trials, with the
    per-trial sequence given as Search and Recall, then Recall Task, then Post-Trial
    Survey.}
  \label{fig:procedure}
\end{figure*}
